\begin{exercise}{19}
Show that the nerve of a category $\mathcal C$ is $2$-coskeletal,
i.e., that the canonical map
$\mathrm N(\mathcal C)\to\mathrm{cosk}_2(\mathrm N(\mathcal C))$
is an isomorphism of simplicial sets.

\begin{proof}
By Definition~1.1.38, the canonical map in degree $m$ is restriction
\[
 \operatorname{Hom}(\Delta^m,\mathrm N(\mathcal C))
 \longrightarrow
 \operatorname{Hom}(\mathrm{sk}_2(\Delta^m),\mathrm N(\mathcal C)).
\]
Let $\varphi$ lie in its target. Its vertices give objects
$c_0,\ldots,c_m$ of $\mathcal C$, and its edges give morphisms
$a_{ij}\colon c_i\to c_j$ for $i<j$. Set
$a_{ii}=\mathrm{id}_{c_i}$, as required by the degenerate edges.
For every $i<j<\ell$, the image of the corresponding triangle is
a functor $[2]\to\mathcal C$, so
\[
 a_{i\ell}=a_{j\ell}a_{ij}.
\]
Together with the identity equations, these are precisely the
composition conditions for a functor $A\colon[m]\to\mathcal C$
sending $i$ to $c_i$ and $i\to j$ to $a_{ij}$.
Associativity is inherited from $\mathcal C$; in particular, for
$i<j<\ell<r$ the two composites of three arrows both equal
$a_{\ell r}a_{j\ell}a_{ij}$.

The $m$-simplex represented by $A$ agrees with $\varphi$ on vertices
and edges, and hence on triangles, since a functor $[2]\to\mathcal C$
is determined by its two successive arrows. These simplices generate
$\mathrm{sk}_2(\Delta^m)$ by Lemma~1.1.39(1), so $A$ extends
$\varphi$. It is unique because its objects and all its arrows
have already been specified. All higher faces and degeneracies are
the restrictions of this single functor, so they are compatible.
Thus restriction is bijective in every degree, including $m\le2$,
and the canonical map is an isomorphism.
\end{proof}
\end{exercise}
